\noindent\textbf{Elección del procedimiento.} En cada pendiente terminal, $Mg\sin\varphi_i=\mu_fMg\cos\varphi_i+k_DV_i^2$. Multiplicar la primera ecuación por $\cos\varphi_2$ y la segunda por $\cos\varphi_1$ elimina el rozamiento al restar. Usar $\sin(\varphi_2-\varphi_1)$ permite aislar $k_D$.

\begin{enumerate}

\item Resolver las dos ecuaciones para la resistencia cuadrática.

\[\fitmath{\mathrm{kD}=\frac{M\cdot g\cdot \sin\left(\varphi_{2}-\varphi_{1}\right)}{{V_{2}}^{2}\cdot \cos\left(\varphi_{1}\right)-{V_{1}}^{2}\cdot \cos\left(\varphi_{2}\right)}}\]

\[\fitmath{\mathrm{kD}=\frac{(80)\cdot (9.81)\cdot \sin\left((0.5235987756)-(0.3490658504)\right)}{{(37.86)}^{2}\cdot \cos\left((0.3490658504)\right)-{(30.43)}^{2}\cdot \cos\left((0.5235987756)\right)}\approx 0.2500489\,\mathrm{kg/m}}\]

\item Sustituir en el primer equilibrio y despejar rozamiento.

\[\fitmath{\mathrm{muf}=\frac{M\cdot g\cdot \sin\left(\varphi_{1}\right)-\mathrm{kD}\cdot {V_{1}}^{2}}{M\cdot g\cdot \cos\left(\varphi_{1}\right)}}\]

\[\fitmath{\mathrm{muf}=\frac{(80)\cdot (9.81)\cdot \sin\left((0.3490658504)\right)-(0.250048886)\cdot {(30.43)}^{2}}{(80)\cdot (9.81)\cdot \cos\left((0.3490658504)\right)}\approx 0.05000322\,}\]

\item El dato de densidad permite convertir $k_D$ al producto área por coeficiente.

\[\fitmath{\mathrm{CDA}=\frac{2\cdot \mathrm{kD}}{\rho}}\]

\[\fitmath{\mathrm{CDA}=\frac{2\cdot (0.250048886)}{(1.2)}\approx 0.4167481\,\mathrm{m^{2}}}\]

\end{enumerate}
